% !PRINTED 535
\weberpage{535}
%<<CONT
$F(x) = 0$, und wenn $\varphi(x) = 0$ gelöst ist, so sind damit alle
Wurzeln von $F(x)$ bekannt. Es genügt also, wenn wir uns auf
die Betrachtung \gesperrt{irreducibler Abel'scher Gleichungen}
beschränken.

\quad Die \gesperrt{Galois'sche} Gruppe einer \gesperrt{Abel'schen} Gleichung, sei
sie irreducibel oder nicht (wenn nur ihre Wurzeln von einander
verschieden sind), hat die Eigenschaft, dass bei der Zusammensetzung
ihrer Permutationen das commutative Gesetz gilt; sie
ist also eine \gesperrt{commutative Gruppe}.

\quad Es seien nämlich
\weq{4}{\alpha,\ \alpha_1 = \Theta_1(\alpha),\ \alpha_2 = \Theta_2(\alpha) \ldots \alpha_{m-1} = \Theta_{m-1}(\alpha)}
die Wurzeln von $F(x)$ und
\weq{5}{\alpha,\ \alpha' = \Theta'(\alpha),\ \alpha'' = \Theta''(\alpha) \ldots}
die darunter enthaltenen Wurzeln des irreducibeln Factors $\varphi(x)$
von $F(x)$.

\quad Da, wie schon bemerkt, $\varphi(x) = 0$ eine \gesperrt{Galois'sche} Resolvente
ist, so besteht die Gruppe der Gleichung aus den Substitutionen
des Körpers $\Omega(\alpha)$, also aus den Substitutionen
\weq{6}{\sigma = (\alpha, \alpha),\ \sigma' = (\alpha, \alpha'),\ \sigma'' = (\alpha, \alpha'') \ldots}

\quad Diese Gruppe befolgt aber das commutative Gesetz; denn
es sei
\[
\begin{aligned}
\sigma' &= (\alpha, \alpha') = [\alpha, \Theta'(\alpha)]\\
\sigma'' &= (\alpha, \alpha'') = [\alpha, \Theta''(\alpha)];
\end{aligned}
\]
dann ist nach \S.\,147
\[
\begin{aligned}
\sigma'\sigma'' &= [\alpha, \Theta'(\alpha)][\Theta'(\alpha), \Theta'\Theta''(\alpha)] = [\alpha, \Theta'\Theta''(\alpha)]\\
\sigma''\sigma' &= [\alpha, \Theta''(\alpha)][\Theta''(\alpha), \Theta''\Theta'(\alpha)] = [\alpha, \Theta''\Theta'(\alpha)].
\end{aligned}
\]
Wegen (3) ist daher $\sigma'\sigma'' = \sigma''\sigma'$, worin $\sigma', \sigma''$ zwei beliebige der
Substitutionen $\sigma$ sein können. Folglich ist die Gruppe der
Substitutionen des Körpers $\Omega(\alpha)$ und damit auch die isomorphe
Permutationsgruppe der Gleichung $F(x) = 0$ \gesperrt{commutativ}.

\quad Es gilt nun auch das Umgekehrte, wobei aber die \gesperrt{Irreducibilität}
vorausgesetzt sein muss.

\quad \gesperrt{Eine irreducible Gleichung $\varphi(x) = 0$ mit commutativer Gruppe ist eine Abel'sche Gleichung.}

\quad Es seien $\alpha, \alpha_1, \alpha_2 \ldots \alpha_{m-1}$ die Wurzeln von $\varphi(x) = 0$ und
$P$ die Gruppe dieser Gleichung, die wegen der Irreducibilität von
$\varphi(x)$ transitiv ist, und die wir jetzt ausserdem als commutativ
annehmen. Es sei $Q$ der Theiler von $P$, der das Element $0$ in
%CONT>>
